\chapter{Prefazione}

\introduction{Mario Tanda}{
    \textit{Il corso fornisce gli strumenti della teoria della probabilità e della statistica che servono a descrivere,
    elaborare e misurare l'informazione: dal calcolo combinatorio alle variabili aleatorie, discrete e continue, e ai
    processi aleatori; dall'inferenza statistica alla teoria della decisione e della stima; fino alla misura
    dell'informazione, con l'entropia e la compressione dei dati.}
  }
\pagestyle{plain}

% ============================================================
\section{Organizzazione del corso}
% ============================================================

Il corso vale 6 CFU (48 ore di lezione), tutte in presenza: circa tre quarti di teoria e un quarto di esercitazioni.
L'esame è una prova scritta integrata da una prova orale. Gli argomenti si dividono in tre blocchi:
\begin{itemize}
  \item \textbf{Teoria della probabilità}: analisi combinatoria, probabilità su spazi finiti, variabili aleatorie
        discrete, variabili aleatorie continue, variabili aleatorie multiple, processi aleatori.
  \item \textbf{Statistica}: statistica inferenziale, teoria della decisione, inferenza bayesiana e non bayesiana,
        stima ricorsiva.
  \item \textbf{Informazione e sua misura} (cenni): entropia, compressione dei dati, mutua informazione, divergenza.
\end{itemize}
Il testo di riferimento, da affiancare alle lezioni, è Sheldon M. Ross, \emph{Probabilità e statistica per
l'ingegneria e le scienze}, Maggioli Editore, IV edizione.

% ============================================================
\section{Come è fatta questa dispensa}
% ============================================================

Questa dispensa segue le slide del corso, una lezione dopo l'altra: ogni parte corrisponde a un pacchetto di slide.
La prima, \emph{Probabilità e variabili aleatorie discrete}, copre il primo; la seconda, \emph{Variabili aleatorie
continue, vettori e processi aleatori}, copre il secondo, fino ai vettori gaussiani e ai processi stazionari; la terza,
\emph{Statistica inferenziale e misura dell'informazione}, copre il terzo: legge dei grandi numeri e teorema centrale del
limite, decisioni e test di ipotesi, stima bayesiana e classica, minimi quadrati, entropia e codifica di sorgente.
L'obiettivo è che ogni formula si possa \textbf{capire}, non solo ricordare: per questo ogni risultato è accompagnato
dal suo perché, da un disegno quando aiuta, e da una dimostrazione scritta passo per passo quando la dimostrazione
esiste. Gli strumenti di matematica che servono (insiemi, sommatorie, limiti, derivate, serie e, nella seconda e nella
terza parte, integrali, integrali doppi, derivate parziali, vettori e matrici, gradienti, autovalori, funzioni convesse)
sono spiegati da zero, nel punto esatto in cui servono, dentro i riquadri viola dei \textbf{richiami di matematica}: chi li conosce già li può saltare, chi non li
ricorda li trova lì, senza dover aprire un altro libro.

% ============================================================
\section{I simboli che incontreremo}
% ============================================================

La matematica usa molti simboli: sono abbreviazioni, e ognuno si legge con parole semplici. Queste tabelle li
raccolgono tutti, divisi per parte; il capitolo indicato è quello in cui il simbolo viene spiegato con calma.

\begin{table}[H]
\centering
\renewcommand{\arraystretch}{1.22}
\begin{tabular}{@{}clP{7.2cm}c@{}}
  \toprule
  \textbf{Simbolo} & \textbf{Si legge} & \textbf{Significato} & \textbf{Cap.} \\
  \midrule
  $x \in A$, $x \notin A$ & $x$ appartiene (non appartiene) ad $A$ & $x$ è (non è) un elemento dell'insieme $A$ & 1 \\
  $A \subseteq B$ & $A$ è contenuto in $B$ & ogni elemento di $A$ sta anche in $B$ & 1 \\
  $\emptyset$, $\card{A}$ & insieme vuoto, cardinalità di $A$ & l'insieme senza elementi; il numero di elementi di $A$
    & 1 \\
  $A \cup B$, $A \cap B$ & $A$ unione $B$, $A$ intersezione $B$ & si verifica almeno uno dei due; si verificano
    entrambi & 1 \\
  $\compl{A}$, $A \setminus B$ & complementare di $A$, $A$ meno $B$ & non si verifica $A$; si verifica $A$ ma non $B$ & 1 \\
  $A \times B$ & $A$ per $B$ (prodotto cartesiano) & l'insieme delle coppie ordinate $(a, b)$ & 1 \\
  $\sum_{i=1}^{n}$, $\prod_{i=1}^{n}$ & sommatoria, produttoria per $i$ da 1 a $n$ & somma, prodotto di $n$ termini
    & 2 \\
  $n!$, $\binom{n}{k}$ & $n$ fattoriale, $n$ su $k$ & $1 \cdot 2 \cdots n$; i sottoinsiemi di $k$ elementi fra $n$ & 2 \\
  $\simeq$ & circa uguale & i due numeri sono vicini & 2 \\
  $a_n \to L$, $\lim\limits_{n \to \infty}$ & $a_n$ tende a $L$, limite per $n$ che tende a infinito & i termini si
    avvicinano a $L$ quanto si vuole & 3 \\
  $P(A)$, $P(A \mid B)$ & probabilità di $A$, di $A$ dato $B$ & quanto è facile che si verifichi $A$ (sapendo che si
    è verificato $B$) & 3 \\
  $\Longrightarrow$, $\iff$ & implica, se e solo se & dalla prima affermazione segue la seconda; le due sono
    equivalenti & 1 \\
  $\forall$ & per ogni & la frase vale per tutti gli elementi indicati & 4 \\
  $X \sim \ldots$ & $X$ si distribuisce come\ldots & $X$ ha la legge indicata a destra, per esempio $\mathcal{B}(N, p)$
    & 5 \\
  $f'(x)$, $\mathrm{d}f / \mathrm{d}x$ & derivata di $f$ & la pendenza del grafico di $f$ nel punto $x$ & 7 \\
  \bottomrule
\end{tabular}
\caption{I simboli della prima parte.}
\end{table}

\begin{table}[H]
\centering
\renewcommand{\arraystretch}{1.22}
\begin{tabular}{@{}clP{7.2cm}c@{}}
  \toprule
  \textbf{Simbolo} & \textbf{Si legge} & \textbf{Significato} & \textbf{Cap.} \\
  \midrule
  $\int_a^b f(x) \dd x$ & integrale di $f$ da $a$ a $b$ & l'area sotto il grafico di $f$ fra $a$ e $b$ & 10 \\
  $f_X(x)$, $F_X(x)$ & densità, CDF di $X$ & probabilità per unità di lunghezza; $P(X \leq x)$ & 10 \\
  $u(x)$, $\mathrm{rect}(x)$ & gradino, rettangolo & 1 per $x \geq 0$ (0 altrove); 1 per $\abs{x} < \frac12$ (0 altrove)
    & 11 \\
  $\iint_C f \dd x \dd y$ & integrale doppio di $f$ su $C$ & il volume sotto la superficie $f$, sopra la regione $C$ & 14 \\
  $\partial F / \partial x$ & derivata parziale rispetto a $x$ & la derivata fatta tenendo fisse le altre
    variabili & 14 \\
  $\mathcal{N}(\mu, \sigma^2)$, $Q(x)$ & gaussiana, funzione $Q$ & la campana di media $\mu$ e varianza $\sigma^2$; la
    sua coda $P(X_0 > x)$ & 15 \\
  $\vett{x}$, $\matr{A}$ & vettore $x$, matrice $A$ & una colonna di numeri; una tabella di numeri & 16 \\
  $\matr{A}^T$, $\matr{A}^{-1}$, $\det\matr{A}$ & trasposta, inversa, determinante & righe e colonne scambiate; la
    matrice che disfa $\matr{A}$; il fattore per cui $\matr{A}$ moltiplica le aree & 16 \\
  $X(n)$ & $X$ di $n$ & un processo aleatorio: un valore aleatorio per ogni istante $n$ & 17 \\
  \bottomrule
\end{tabular}
\caption{I simboli della seconda parte.}
\end{table}

\begin{table}[H]
\centering
\renewcommand{\arraystretch}{1.22}
\begin{tabular}{@{}clP{7.2cm}c@{}}
  \toprule
  \textbf{Simbolo} & \textbf{Si legge} & \textbf{Significato} & \textbf{Cap.} \\
  \midrule
  $\overline{X}_n$ & media campione & la media dei dati $X_1, \ldots, X_n$ & 18 \\
  i.i.d. & indipendenti e identicamente distribuite & dati che sono prove ripetute dello stesso esperimento & 18 \\
  $\Phi(x)$ & fi di $x$ & la CDF della gaussiana standard, $1 - Q(x)$ & 18 \\
  $\decidi{H_1}{H_2}$ & maggiore: $H_1$; minore: $H_2$ & se vale $>$ si decide $H_1$, se vale $<$ si decide $H_2$ & 19 \\
  $L(\vett{x})$ & rapporto di verosimiglianza & quanto i dati sono più probabili sotto un'ipotesi che sotto l'altra
    & 19 \\
  $\lfloor a \rfloor$, $\lceil a \rceil$ & parte intera inferiore, superiore & $a$ arrotondato per difetto, per eccesso
    all'intero & 19, 26 \\
  $\alpha$, $\beta$ & falso allarme, mancata rivelazione & le probabilità di errore di prima e di seconda specie;
    $1 - \beta$ è la potenza & 20 \\
  $Q^{-1}(\alpha)$ & Q inversa di alfa & il numero $x$ per cui $Q(x) = \alpha$ & 20 \\
  $\hat\theta$, $\hat\Theta$ & theta stimato & una stima; uno stimatore & 21 \\
  $\arg\min$, $\arg\max$ & argomento del minimo, del massimo & il punto in cui una funzione è minima, massima & 21 \\
  $\Gamma(x)$ & gamma di $x$ & il fattoriale esteso: $\Gamma(n + 1) = n!$ & 21 \\
  $b_n(\theta)$, $I_n(\theta)$ & polarizzazione, informazione di Fisher & l'errore medio di uno stimatore; quanto i dati
    dicono sul parametro & 22 \\
  $\nabla_{\vett{w}}\, g$ & gradiente di $g$ & il vettore delle derivate parziali & 23 \\
  $\norm{\vett{v}}$ & norma di $\vett{v}$ & la lunghezza del vettore, $\sqrt{\vett{v}^T\vett{v}}$ & 24 \\
  $i(A)$, $H(X)$ & informazione, entropia & $\log \frac{1}{P(A)}$; l'informazione media di $X$ & 25 \\
  \bottomrule
\end{tabular}
\caption{I simboli della terza parte.}
\end{table}

% ============================================================
\section{Convenzioni tipografiche}
% ============================================================

Nel testo si usano alcuni riquadri ricorrenti, ognuno con un colore e un compito precisi.

\dfn{Definizione}{
  I riquadri rossi contengono le \textbf{definizioni}: il significato esatto di una parola nuova. Sono da saper
  ripetere con precisione.
}

\teorema{Teorema}{
  I riquadri blu contengono i \textbf{teoremi}: affermazioni che si possono dimostrare a partire dalle definizioni.
  Sono numerati capitolo per capitolo (Teorema 3.2 è il secondo teorema del capitolo 3), così si possono citare.
}

\begin{dimostrazione}
\passo{Che cosa contiene} Sotto un teorema, quando la dimostrazione esiste, c'è questo riquadro grigio: la
dimostrazione scritta \textbf{un passaggio alla volta}, ciascuno con il suo motivo.
\passo{Come finisce} Il quadratino nero in fondo a destra vuol dire «come volevasi dimostrare»: da lì in poi il
teorema si può usare senza dubbi.
\end{dimostrazione}

\begin{richiamo}{il titolo dice di che cosa}
I riquadri viola spiegano \textbf{da zero} gli strumenti di matematica che servono in quel punto: che cos'è un
limite, una derivata, un integrale, una sommatoria. Sono scritti per chi non li ha mai visti o li ha dimenticati.
\end{richiamo}

\begin{formula}[nome della formula]
\[ \text{le formule da tenere a mente, in evidenza} \]
\end{formula}

\info{Nota}{
  I riquadri verdi contengono \textbf{intuizioni}, analogie e osservazioni che aiutano a capire, ma non sono
  indispensabili.
}

\warning{Attenzione}{
  I riquadri gialli segnalano gli \textbf{errori tipici} e i punti in cui è facile confondersi.
}

\begin{esempio}{Struttura di un esempio}
Gli esempi sono racchiusi fra due righe sottili e numerati all'interno di ogni capitolo.
\end{esempio}

\begin{esercizio}{Struttura di un esercizio}
Gli esercizi sono racchiusi fra righe doppie.
\svolgimento
Lo svolgimento viene subito dopo il testo: conviene provarci da soli prima di leggerlo.
\end{esercizio}
